% ============================================================
%  Self-contained standalone paper on the width-three slices c=3 and c=4.
%  First xelatex/pdflatex run writes heffter_c3.bib via filecontents;
%  then bibtex -> (pdf|xe)latex -> (pdf|xe)latex.  No external .bib needed.
% ============================================================
\begin{filecontents*}[overwrite]{heffter_c3.bib}
@article{Arch2015,
  author  = {Archdeacon, D. S.},
  title   = {Heffter arrays and biembedding graphs on surfaces},
  journal = {Electron. J. Combin.},
  volume  = {22},
  year    = {2015},
  note    = {Paper {\#}P1.74; DOI: 10.37236/4874}
}
@article{PT2025,
  author  = {Pellegrini, Marco Antonio and Traetta, Tommaso},
  title   = {Toward a solution of {Archdeacon}'s conjecture on integer {Heffter} arrays},
  journal = {J. Combin. Des.},
  volume  = {33},
  number  = {8},
  pages   = {310--323},
  year    = {2025},
  note    = {DOI: 10.1002/jcd.21983}
}
@misc{MP2025,
  author       = {Morini, Fiorenza and Pellegrini, Marco Antonio},
  title        = {An update on the existence of integer {Heffter} arrays},
  howpublished = {Preprint, \texttt{arXiv:2510.08302}},
  year         = {2025}
}
@article{ABD2017,
  author  = {Archdeacon, Dan S. and Boothby, Tom and Dinitz, Jeffrey H.},
  title   = {Tight {Heffter} arrays exist for all possible values},
  journal = {J. Combin. Des.},
  volume  = {25},
  pages   = {5--35},
  year    = {2017},
  note    = {DOI: 10.1002/jcd.21520}
}
@article{Skolem1957,
  author  = {Skolem, Thoralf},
  title   = {On certain distributions of integers in pairs with given differences},
  journal = {Math. Scand.},
  volume  = {5},
  pages   = {57--68},
  year    = {1957}
}
@article{Zhou2018DM,
  author  = {Zhou, Chen and Li, Wen and Zhang, Yong and Su, Renwang},
  title   = {Existence of centre-complementary magic rectangles},
  journal = {Discrete Math.},
  volume  = {341},
  pages   = {1952--1958},
  year    = {2018},
  note    = {DOI: 10.1016/j.disc.2018.03.022}
}
@article{Zhou2018JCD,
  author  = {Zhou, Chen and Li, Wen and Zhang, Yong and Su, Renwang},
  title   = {Existence of magic 3-dimensional rectangles},
  journal = {J. Combin. Des.},
  volume  = {26},
  pages   = {280--309},
  year    = {2018},
  note    = {DOI: 10.1002/jcd.21593}
}
@incollection{PasottiDinitz2024,
  author={Pasotti, Anita and Dinitz, Jeffrey H.},
  title={A survey of {Heffter} arrays},
  booktitle={New Advances in Designs, Codes and Cryptography},
  editor={Colbourn, Charles J. and Dinitz, Jeffrey H.},
  series={Fields Institute Communications}, volume={86},
  pages={353--392}, publisher={Springer}, year={2024},
  doi={10.1007/978-3-031-48679-1_20}
}
@article{MoriniPellegrini2022,
  author={Morini, Fiorenza and Pellegrini, Marco Antonio},
  title={Rectangular {Heffter} arrays: a reduction theorem},
  journal={Discrete Math.}, volume={345}, number={12}, pages={113073},
  year={2022}, doi={10.1016/j.disc.2022.113073}
}
@article{MoriniPellegrini2020,
  author={Morini, Fiorenza and Pellegrini, Marco Antonio},
  title={On the existence of integer relative {Heffter} arrays},
  journal={Discrete Math.}, volume={343}, number={11}, pages={112088},
  year={2020}, doi={10.1016/j.disc.2020.112088}
}
@article{CostaMoriniPasottiPellegrini2018,
  author={Costa, Simone and Morini, Fiorenza and Pasotti, Anita and Pellegrini, Marco Antonio},
  title={A problem on partial sums in abelian groups},
  journal={Discrete Mathematics}, volume={341}, number={3}, pages={705--712},
  year={2018}, doi={10.1016/j.disc.2017.11.013}
}
@article{BurattiPasotti2024,
  author={Buratti, Marco and Pasotti, Anita},
  title={Heffter spaces}, journal={Finite Fields Appl.}, volume={98}, pages={102464},
  year={2024}, doi={10.1016/j.ffa.2024.102464}
}
@article{BurattiPasotti2025,
  author={Buratti, Marco and Pasotti, Anita},
  title={More Heffter spaces via finite fields}, journal={J. Combin. Des.},
  volume={33}, pages={188--194}, year={2025}, doi={10.1002/jcd.21974}
}
\end{filecontents*}

\documentclass[11pt]{article}
\usepackage{amsmath,amssymb,amsthm}
\usepackage[a4paper,margin=2.5cm]{geometry}
\usepackage{xcolor}
\usepackage[colorlinks=true,linkcolor=blue!55!black,urlcolor=teal!60!black,citecolor=blue!55!black]{hyperref}
\usepackage{bookmark}
\usepackage{enumitem}
\usepackage{longtable}

\theoremstyle{plain}
\newtheorem{theorem}{Theorem}
\newtheorem{proposition}{Proposition}
\newtheorem{corollary}{Corollary}
\newtheorem{lemma}{Lemma}
\theoremstyle{definition}
\newtheorem{definition}{Definition}
\newtheorem{problem}{Problem}

\newcommand{\Z}{\mathbb{Z}}
\newcommand{\Hm}{\mathrm{H}}
\newcommand{\IHS}{\mathrm{IHS}}

\title{\textbf{Existence of integer Heffter array sets $\IHS(m,3;3)$ and $\IHS(m,3;4)$}}
\author{%
\begin{tabular}{@{}c@{\qquad\qquad}c@{}}
\textbf{Yong Zhang}\textsuperscript{*} & \textbf{Guangzhou Chen}\\[2pt]
\begin{tabular}[t]{@{}c@{}}
School of Mathematics and Statistics\\
Yancheng Teachers University\\
Yancheng 224002, P.~R.~China\\
\texttt{zyyctc@126.com}
\end{tabular}
&
\begin{tabular}[t]{@{}c@{}}
School of Mathematics and Statistics\\
Henan Normal University\\
Xinxiang 453007, P.~R.~China\\
\texttt{chenguangzhou0808@163.com}
\end{tabular}
\\[4pt]
\multicolumn{2}{c}{\footnotesize $^*$ Corresponding author}
\end{tabular}}
\date{}

\begin{document}
\maketitle

\begin{abstract}
We completely determine two fixed slices of the width-three integer Heffter array set problem.
An $\IHS(m,3;3)$ exists for every $m\equiv3\pmod4$ with $\gcd(m,3)=1$, and an
$\IHS(m,3;4)$ exists for every odd $m\ge5$ with $\gcd(m,3)=1$. The construction starts from
a closed formula for centre-complementary $3\times(2n+1)$ Kotzig arrays. A bounded lift of
their positive half-columns supplies zero-sum rows; periodic offset words correct the support,
and finite boundary and orientation certificates enforce exact coverage and zero column sums.
The bulk rules have periods $24$ and $48$ for $c=3$ and $c=4$, respectively. The finitely many
preperiod parameters are handled by explicit certificates included in the manuscript and the
accompanying repository. Thus the infinite families reduce to finite affine endpoint and
signed-sum checks; the verification programs only decode and check the printed data. No claim
is made here for $c\ge5$.
\end{abstract}

\noindent\textbf{Keywords:} integer Heffter array; Heffter array set; Skolem sequence; Kotzig array;
finite verification.\\
\textbf{2020 Mathematics Subject Classification:} 05B20, 05B30.

\section{Introduction}\label{sec:intro}

Heffter arrays were introduced by Archdeacon~\cite{Arch2015} in the study of orientable
biembeddings of complete graphs. Their definitions and principal existence results are surveyed
by Pasotti and Dinitz~\cite{PasottiDinitz2024}. Related work concerns partial sums in abelian
groups~\cite{CostaMoriniPasottiPellegrini2018}, integer relative arrays~\cite{MoriniPellegrini2020},
rectangular arrays~\cite{MoriniPellegrini2022}, and the recently introduced Heffter spaces
~\cite{BurattiPasotti2024,BurattiPasotti2025}.

\begin{definition}[integer Heffter array \cite{Arch2015}]
An integer Heffter array $\Hm(m,n;s,k)$ is an $m\times n$ partially filled array, with entries from
$\{\pm1,\pm2,\dots,\pm nk\}\subset\Z$, such that each row has exactly $s$ filled cells and each column
exactly $k$; the absolute values of the entries are pairwise distinct; and every row sum and column sum
is $0$.
\end{definition}

For an integer Heffter array $\Hm(m,n;s,k)$, the necessary conditions are
$ms=nk$, $3\le s\le n$, $3\le k\le m$, and $nk\equiv0,3\pmod4$; Archdeacon's conjecture
asserts that these conditions are sufficient. In the coprime odd-width regime, Pellegrini and
Traetta~\cite{PT2025} introduced integer Heffter array sets, and Morini and Pellegrini~\cite{MP2025}
reduced the remaining cases to two short-width problems.

\begin{definition}[integer Heffter array set \cite{PT2025}]
An $\IHS(m,n;c)$ is a set of $c$ arrays of size $m\times n$ whose union has $mnc$ entries with
pairwise distinct absolute values, collectively equal to $[1,mnc]$. Every row and every column of each
array has sum zero.
\end{definition}

When $s$ and $k$ are coprime, write $m=ck$ and $n=cs$. In this notation, an $\IHS(k,s;c)$
assembles an $\Hm(ck,cs;s,k)$ \cite[proof of Theorem~C]{MP2025}. The remaining short-width
cases therefore have $k=3$ or $k=5$. We study the width-three case:

\begin{problem}[{\cite[Problem~1]{MP2025}}]\label{prob:1}
For every $c\ge1$ and every odd $m\ge3$ with $\gcd(m,3)=1$ and $mc\equiv0,1\pmod4$, construct an
$\IHS(m,3;c)$.
\end{problem}

The parameter $c$ separates Problem~\ref{prob:1} into fixed slices. The case $c=1$ is already
settled: an $\IHS(m,3;1)$ is precisely a
\emph{tight} integer Heffter array $\Hm(m,3;3,m)$---an $m\times3$ array with no empty cells whose
support is exactly $[1,3m]$---and Archdeacon, Boothby and Dinitz proved that tight Heffter arrays
exist for all admissible parameters~\cite{ABD2017}; for width three this covers every admissible $m$
(here $m\equiv1\pmod4$, $\gcd(m,3)=1$). The case $c=2$ is vacuous in Problem~\ref{prob:1}, since its
odd parameter $m$ gives $mc=2m\equiv2\pmod4$. Thus the first nonvacuous case beyond one array is
$c=3$, admissible exactly when $3m\equiv1\pmod4$, i.e.\ $m\equiv3\pmod4$. Our first main result
settles it. The next even slice, $c=4$, is admissible for every odd $m$ coprime to $3$; our second
main result settles that slice as well.

\begin{theorem}\label{thm:main}
For every integer $m\ge3$ with $m\equiv3\pmod4$ and $\gcd(m,3)=1$, an $\IHS(m,3;3)$ exists.
\end{theorem}

Together with the tight case $c=1$ of~\cite{ABD2017}, Theorem~\ref{thm:main} settles the $c=1$
and $c=3$ slices of Problem~\ref{prob:1}. Its construction is based on the closed formula for
the centre-complementary Kotzig arrays of Zhou, Li, Zhang and Su~\cite{Zhou2018DM}. Theorem
~\ref{thm:c4full} below settles the next admissible slice, $c=4$. Both proofs use the same
reduction into a support partition, a grouping step, and an orientation certificate, but the
periodic data and boundary states differ between the two slices.

The results are fixed-slice theorems, not a solution for arbitrary $c$. In particular, no claim is
made here for $c\ge5$; the concluding section describes which parts of the certificate method may
be useful for those cases. The displayed finite data, together with the three certificate source
files listed in the data-availability paragraph, make the two constructions independently checkable.

\paragraph{Notation.}
$[1,N]=\{1,\dots,N\}$. A \emph{sum-triple} is a set $\{x,y,x+y\}$ of three distinct positive integers.
A multiset is \emph{signable} if some choice of signs makes it sum to $0$ (equivalently, its total is
even and half of it is a subset sum).

\section{Reduction to sum-triples and balanceability}\label{sec:reduction}

\begin{proposition}[width-three reduction]\label{prop:red}
Let $m\ge3$, $c\ge1$, $N=mc$. An $\IHS(m,3;c)$ exists if and only if
\textup{(i)} $[1,3N]$ can be partitioned into $N$ sum-triples, and
\textup{(ii)} these $N$ triples can be split into $c$ groups of $m$, each \emph{orientable} (for each
triple choose a sign and a placement of its three signed values into three columns) so that the three
column sums of the group vanish.
\end{proposition}
\begin{proof}
Because no entry is zero and the absolute values are distinct, three signed entries in a zero-sum
row can be reordered as $x$, $y$, and $-(x+y)$, with $x,y>0$. Thus the absolute values in every
row form a sum-triple. Conversely, every sum-triple $\{x,y,x+y\}$ gives a zero-sum row by using
the signs $(x,y,-(x+y))$. The zero column sums are exactly the orientation condition in (ii), while
the $c$ arrays together cover $[1,3N]$ precisely when the triples form the partition in (i).
\end{proof}

\begin{proposition}[Skolem partition]\label{prop:skolem}
The interval $[1,3N]$ can be partitioned into $N$ sum-triples if and only if
$N\equiv0,1\pmod4$.
\end{proposition}
\begin{proof}
Every sum-triple has even total, so necessity is equivalent to the evenness of
$3N(3N+1)/2$, which is exactly $N\equiv0,1\pmod4$. For sufficiency, a Skolem sequence of
order $N$ exists precisely when $N\equiv0,1\pmod4$~\cite{Skolem1957}. Hence there are pairs
$\{x_i,y_i\}$, $i\in[1,N]$, partitioning $[1,2N]$ such that $y_i-x_i=i$. The sets
$\{i,N+x_i,N+y_i\}$ partition $[1,3N]$ and are sum-triples because
$i+(N+x_i)=N+y_i$.
\end{proof}

Thus the support partition exists under exactly the necessary congruence condition. It remains to
balance the triples within each group. Let $C_1,C_2,C_3$ denote the three column sums of a fixed
group. Since every row already sums to zero, $C_1+C_2+C_3=0$; hence zero column sums are
equivalent to the two conditions $C_1+C_2=0$ and $C_1-C_2=0$. For a triple $\{a,b,c\}$ with
$c=a+b$ and $a<b$, a legal orientation is a \emph{mode} $\mathrm m\in\{a,b,c\}$ (which element
occupies the third column), together with an overall sign and a possible interchange of the first
two columns. Writing $\mu$ for the magnitude of its contribution to $C_1{+}C_2$ and $|w|$ for
that to $C_1{-}C_2$,
\[
(\mathrm m{=}c):\ \mu=c,\ |w|=b-a;\quad
(\mathrm m{=}a):\ \mu=a,\ |w|=b+c;\quad
(\mathrm m{=}b):\ \mu=b,\ |w|=a+c.
\]

\begin{lemma}[exact characterisation of balancing]\label{lem:exact}
A group of sum-triples can be oriented to zero column sums if and only if it has a mode assignment
$(\mathrm m_i)_i$ for which both multisets $\{\mu_i\}_i$ and $\{|w_i|\}_i$ are signable.
\end{lemma}
\begin{proof}
For a sign $\varepsilon\in\{\pm1\}$, the three modes may be represented by the signed entries
\[
\begin{array}{c|c|c}
\text{mode}&\text{first two columns}&\text{third column}\\\hline
c&(\varepsilon a,\varepsilon b)&-\varepsilon c\\
a&(\varepsilon c,-\varepsilon b)&-\varepsilon a\\
b&(\varepsilon c,-\varepsilon a)&-\varepsilon b.
\end{array}
\]
Interchanging the first two columns independently changes the sign of their difference. Hence a row
contributes $\varepsilon\mu$ to $C_1+C_2$ and an independently signed value of magnitude $|w|$ to
$C_1-C_2$, with the three pairs $(\mu,|w|)$ displayed above. The overall signs therefore realize
any signing of the $\mu_i$, while the column interchanges independently realize any signing of the
$|w_i|$. Thus $C_1=C_2=0$ if and only if both associated multisets are signable. These modes,
overall signs, and interchanges exhaust all zero-sum orientations of the rows.
\end{proof}

\begin{corollary}[balance certificate]\label{cor:cert}
A group is balanceable if and only if it admits a mode assignment $(\mathrm m_i)_i$ and subsets $U,V$ of indices
with $\sum_{i\in U}\mu_i=\sum_{i\notin U}\mu_i$ and $\sum_{i\in V}|w_i|=\sum_{i\notin V}|w_i|$; the
orientation is then read off constructively.
\end{corollary}

\section{Kotzig lift and periodic support correction}\label{sec:construction}

\subsection{The half-column lift}

Let $n=3m$ and $r=2n+1$. Zhou--Li--Zhang--Su~\cite[Theorem~2.2]{Zhou2018DM} give the $3\times r$
centre-complementary Kotzig array $A=(a_{i,j})$, $i\in[-1,1]$, $j\in[-n,n]$, by the closed formula
\begin{equation}\label{eq:kotzig}
 a_{i,j}=\frac{1-(-1)^{\,j+n}\,r}{4}\,i+\frac{2-3|i|}{2}\,j ,
\end{equation}
For odd $n$, this formula is integer-valued; each row is a permutation of $[-n,n]$, each column sums
to zero, and $a_{-i,-j}=-a_{i,j}$. For
$j\in[1,n]$,
\begin{equation}\label{eq:vals}
(a_{-1,j},a_{1,j})=
\begin{cases}
\bigl(\tfrac{n-j}{2},\ -\tfrac{n+j}{2}\bigr), & j\equiv n\pmod2,\\[3pt]
\bigl(-\tfrac{n+1+j}{2},\ \tfrac{n+1-j}{2}\bigr), & j\not\equiv n\pmod2 .
\end{cases}
\end{equation}

\begin{proposition}[half-column lift]\label{prop:lift}
Let $n$ be odd. For $j\in[1,n]$ set
$R_j=\bigl(3a_{-1,j}-1,\ 3j,\ 3a_{1,j}+1\bigr)$. Then each $R_j$ is a zero-sum row whose absolute
values form a sum-triple, the $3n$ absolute values are pairwise distinct, and their union is
\[
 [1,3n+1]\setminus\{h_n\},\qquad h_n=\frac{3n+5}{2}.
\]
\end{proposition}
\begin{proof}
Write $n=2\kappa-1$. Row sums vanish because each column of $A$ sums to zero and
$-1+0+1=0$. For an odd index $j=2\ell-1$, $1\le\ell\le\kappa$, put
$p=\kappa-\ell$ and $q=\kappa+\ell-1$. By \eqref{eq:vals}, the absolute values are
$\{|3p-1|,3j,3q-1\}$. When $p=0$, one has $q=j=n$ and $1+(3q-1)=3j$; otherwise
$(3p-1)+3j=3q-1$. These rows supply $1$ and all positive values congruent to $2$ modulo $3$ in
$[1,3n+1]$. For an even index $j=2\ell$, $1\le\ell\le\kappa-1$, put
$s=\kappa-\ell$ and $w=\kappa+\ell$. The absolute values are
$\{3s+1,3j,3w+1\}$ and $(3s+1)+3j=3w+1$; they supply all values congruent to $1$ modulo $3$ except
$3\kappa+1$. Finally, the middle entries supply every positive multiple of $3$ through $3n$.
Consequently all values are distinct and their union is
$[1,3n+1]\setminus\{3\kappa+1\}$, where $3\kappa+1=(3n+5)/2$.
\end{proof}

The lift omits $h_n$ and uses the overshoot $3n+1$. To see why a local repair is insufficient, call
$\{3g-1,3g,3g+1\}$ the $g$-th block. Proposition~\ref{prop:lift} supplies three base cells of
magnitude $g$ for every $g\in[1,n]$ except $g=\kappa=(n+1)/2$, where one base cell would have
magnitude $0$. Thus the missing cell in the $\kappa$-th block cannot be repaired using only the
offsets $(-1,0,1)$. We therefore use zero-sum offsets in $\{-2,\dots,2\}^3$, allowing a bounded
cascade between adjacent blocks while preserving every row sum.

\subsection{The uniform correction rule}\label{sec:rule}

\begin{lemma}[deterministic correction]\label{lem:rule}
Let $n=3m\ge369$, where $m\equiv3\pmod4$ and $\gcd(m,3)=1$. Define an explicit offset assignment
$\delta:[1,n]\to\{-2,\dots,2\}^3$, with coordinate sum zero at every index, by the following
period-$24$ rule. For the class
$n\equiv9\pmod{24}$ there are two $24$-letter offset words $W_1,W_2$ and constant phases $(6,12)$ with
\[
\delta(j)=
\begin{cases}
W_1[(j-6)\bmod24], & 21\le j\le\kappa-21,\\
W_2[(j-12)\bmod24], & \kappa+21\le j\le n-21,
\end{cases}\qquad \kappa=\tfrac{n+1}2,
\]
completed by three bounded boundary tables (for $j\le20$, $|j-\kappa|\le20$, $j\ge n-20$), one set per
residue of $n$ modulo $48$.  The class $n\equiv21\pmod{24}$ uses phases $(10,13)$ and the rotated
words $\widehat W_1[p]=W_1[p+16\bmod24]$, $\widehat W_2[p]=W_2[p+13\bmod24]$.
The two master words and all four boundary-table sets are printed in
Section~\ref{sec:offset-data}. For every admissible $n\ge369$, the resulting corrected rows partition
$[1,3n]$ into $n$ sum-triples.
\end{lemma}
\begin{proof}
Write $n=e+24t$, where $e\in\{9,21\}$. Choose the unique representative
$t=t_0+12Q$ with $15\le t_0\le26$ and $Q\ge0$. Since $n=3m$ and $\gcd(m,3)=1$, exclude
$t_0\equiv0\pmod3$ for $e=9$ and $t_0\equiv1\pmod3$ for $e=21$.  This leaves $16$ admissible
states.  In a bulk word write
$j=j_0+24h$.  In a fixed state the sign of each coordinate of the corrected lift is fixed, so its
absolute value is affine in $(Q,h)$; its $h$-coefficient has magnitude $36$ or $72$.  Split a
magnitude-$36$ family according to the parity of $h$.  Every resulting family is then an arithmetic
progression with difference $72$.

Fix a residue $u\in[1,72]$ and consider the values congruent to $u$ modulo $72$, after subtracting
$u$ and dividing by $72$. The bulk progressions and the boundary singletons become integer intervals
whose endpoints are affine in $Q$. Substitution of the two word pairs and four boundary-table sets printed in
Section~\ref{sec:offset-data} gives, in each admissible state and for every $u$, consecutive
intervals beginning at $0$ and ending at
\[
  12Q+\left\lfloor\frac{3n_0-u}{72}\right\rfloor,
  \qquad n_0=e+24t_0.
\]
The same substitution verifies, in each row, that the three absolute values satisfy
$a+b=c$ after sorting. The endpoint offsets at $Q=0$ are given by the displayed boundary strings,
and the affine endpoint slopes are fixed by the period words; hence the same adjacency holds for every
$Q\ge0$. Thus every integer in $[1,3n]$ occurs once. The calculation consists only of comparing affine
endpoint pairs; the printed words and tables are the complete finite input, so the calculation may be
checked directly without the generating program.
\end{proof}

The rule is finite-state in $m$.  Inside either bulk interval, $\delta(j+24)=\delta(j)$, while the
three boundary tables depend only on $n\bmod48$.  More sharply, equation~\eqref{eq:c3-word-shift}
shows that changing class from $n\equiv9$ to $21$, or conversely, is a half-period shift by $12$ in
both effective bulk layers.  Under $m\mapsto m+4$ the boundary state follows
$9\to21\to33\to45\to9\pmod{48}$.  Thus parameter growth inserts periodic bulk blocks and updates one
of four boundary states; it is not an append-only operation on the arrays.

\section{Grouping, orientation, and the main theorem}\label{sec:proof}

Fix the corrected partition of Lemma~\ref{lem:rule}, with triple $T_j$ indexed by column $j\in[1,n]$,
$n=3m$. Group by residue: $G_r=\{T_j:j\equiv r\pmod3\}$, $r=0,1,2$, each of size $m$.

\begin{lemma}[periodic orientation]\label{lem:orient}
Each group $G_r$ is balanceable, and a balancing is given by an explicit rule: orient $T_j$ in the
decoupled form of Corollary~\ref{cor:cert} (the sum $c_j$ in the third column with sign
$\varepsilon_j$, the two summands in columns $1,2$ in an order set by a sign $\eta_j$), with one
designated row in the $\kappa$-window possibly switched to mode $a$ or $b$ to correct parity. The signs
$(\varepsilon_j,\eta_j)$ are constant on the three bulk zones determined by the bounded windows at
$m$ and $\kappa$, after further subdivision by $j\bmod24$; the four boundary windows are at
$1,m,\kappa,n$. There are $24$ such sign patterns---one for each of the eight admissible residues of
$n$ modulo $144$ and each of the three groups. Every pattern
fixes the signs of both chains $\{c_j\}$ and $\{b_j-a_j\}$ and is printed in
Section~\ref{sec:orientation-data}.
\end{lemma}
\begin{proof}
By Lemma~\ref{lem:exact} it suffices to exhibit, per group, a mode assignment making both chains
signable; the decoupled certificate of Corollary~\ref{cor:cert} does this with all modes equal to $c$
except the one parity-switching row.  The complete $24$ patterns are printed in
Section~\ref{sec:orientation-data}. To certify the infinite tail, write
$n=e+24(\tau+12Q)$ with $e\in\{9,21\}$, $15\le\tau\le26$, and $Q\ge0$.  The coprimality restriction
leaves $16\cdot3=48$ admissible group states. The stored pattern indexed by
$s=15+((\tau-15)\bmod6)$ is used for both $\tau=s$ and $\tau=s+6$. In a fixed state and sign
category, each contribution to either chain is a sum of affine values over an arithmetic progression
with affine endpoints, hence a polynomial in $Q$ of degree at most two. Direct substitution of the
decoded signs gives zero at $Q=0,1,\ldots,6$ for both chains and for both values of $\tau$ represented
by each stored row. Each signed total is therefore the zero polynomial. This is a finite arithmetic
certificate, not an assumption about the output of a search program.
\end{proof}

\begin{lemma}[finite bridge]\label{lem:bridge}
The same offset and orientation rules produce an $\IHS(m,3;3)$ for
\[
\mathcal B=\{67,71,79,83,91,95,103,107,115,119\}.
\]
\end{lemma}
\begin{proof}
For $m\ge67$ the four orientation windows at $1,m,\kappa,n$ are disjoint; in particular,
$\kappa-m=(m+1)/2\ge34>12+20$. Hence the category rule in
Section~\ref{sec:orientation-data} applies without overlap. For every $m\in\mathcal B$, substitute
$n=3m$ in the offset words and select the stored orientation row prescribed there. The ten resulting
checks are printed in Section~\ref{sec:bridge-data}: the absolute-value set has minimum $1$, maximum
$9m$, and cardinality $9m$, while all three column-sum vectors are zero. Row sums are zero by the
construction, so these data prove the assertion.
\end{proof}

\begin{proposition}[alternating transport of the canonical form]\label{prop:mtransport}
Let $m\ge67$ be admissible and put
\[
 d(m)=\begin{cases}4,&m\equiv1\pmod3,\\8,&m\equiv2\pmod3.\end{cases}
\]
Write $R_j^{(m)}$ and $\delta_j^{(m)}$ for the corrected row and its offset in the construction above.
For every inherited index $1\le j\le3m$,
\begin{equation}\label{eq:c3-standalone-transport}
 R_j^{(m+d)}=R_j^{(m)}+\frac{9d}{2}(-1)^{j+3m}(1,0,-1)
 +\bigl(\delta_j^{(m+d)}-\delta_j^{(m)}\bigr).
\end{equation}
The remaining $3d$ target rows are new, with $d$ in each group modulo $3$.  The support and
orientation states are the next states in the printed finite tables, so the construction follows the
single chain
\[
 67\xrightarrow{+4}71\xrightarrow{+8}79\xrightarrow{+4}83
 \xrightarrow{+8}91\xrightarrow{+4}\cdots .
\]
This is a recurrence between the canonical corrected-Kotzig forms, not an extension operator on an
arbitrary presentation of an $\IHS(m,3;3)$.
\end{proposition}
\begin{proof}
Subtract the closed Kotzig formula at $3m$ from the formula at $3(m+d)$ and then add the two offset
vectors; this gives~\eqref{eq:c3-standalone-transport}.  The count of new rows is immediate modulo
$3$.  Equation~\eqref{eq:c3-word-shift} gives the half-period bulk update for the $+4$ step, the four
boundary tables give the stated boundary-state cycle, and Lemmas~\ref{lem:rule} and~\ref{lem:orient}
certify the target support and balance states.  No new search is used.
\end{proof}

An append-only interpretation is impossible.  If the old entries were left fixed in a $+4$ step,
each new zero-sum row would have to use values from $[9m+1,9m+36]$; but
$(9m+1)+(9m+2)>9m+36$ for $m\ge4$, so no three values in that interval form a sum-triple.

\begin{proof}[Proof of Theorem~\ref{thm:main}]
Let $m\equiv3\pmod4$ with $\gcd(m,3)=1$ and put $n=3m$. Then $n\equiv1\pmod4$ and $n$ is odd. For
$m<67$, the ten admissible values are
\[
 7,11,19,23,31,35,43,47,55,59.
\]
Complete signed arrays are displayed below for $m=7$, while
Section~\ref{sec:small-data} gives three exceptional support words for $m=11,19,23$ and six
preperiod orientation certificates for $m=31,35,43,47,55,59$; the latter six already use the uniform
support rule.
Lemma~\ref{lem:bridge} handles the ten admissible values from $67$ to $119$. For every remaining
admissible $m>119$ (the first is $127$),
Lemma~\ref{lem:rule} gives the exact sum-triple partition of $[1,3n]$; grouping by $j\bmod3$ and
 orienting by Lemma~\ref{lem:orient} gives three zero-sum $m\times3$ arrays whose absolute-value
supports partition $[1,9m]$.
\end{proof}

\begin{center}
\emph{Example} ($m=7$). The smallest admissible instance is the following $\IHS(7,3;3)$:
{\footnotesize
\[
\setlength{\arraycolsep}{3pt}
\left[\begin{array}{rrr}
-31&32&-1\\10&-39&29\\-21&41&-20\\19&-47&28\\37&12&-49\\-11&57&-46\\-3&-56&59
\end{array}\right]\quad
\left[\begin{array}{rrr}
36&-30&-6\\-38&25&13\\22&23&-45\\-48&33&15\\40&-54&14\\50&-55&5\\-62&58&4
\end{array}\right]\quad
\left[\begin{array}{rrr}
34&-27&-7\\42&-24&-18\\-43&26&17\\-51&16&35\\9&-53&44\\-52&60&-8\\61&2&-63
\end{array}\right]
\]
}
\end{center}
\noindent
Each row sums to $0$, each array's three column sums are $0$, and the union of the absolute-value
supports of the three arrays is exactly $[1,63]$.

\section{Explicit finite data for the complete \texorpdfstring{$c=3$}{c=3} slice}\label{sec:certificates}
This section prints every finite witness used by the periodic construction and the finite bridge.
Thus those parts of the existence proof are logically independent of the accompanying programs:
the programs merely regenerate and check the data below.  All strings are read from left to right.
Throughout this section \(n=3m\), \(\kappa=(n+1)/2\), and all phases are taken modulo \(24\).
The symbols \(L\), \(M\), and \(U\) in the boundary table refer to the lower endpoint,
the \(\kappa\)-window, and the upper endpoint, respectively.  The bulk blocks are denoted
by \(B_1,B_2,B_3\).  Thus a certificate is read in three steps: decode the offset word,
lift it with the formula in Section~\ref{sec:construction}, and then apply the displayed orientation
and sign strings.

\subsection{Offset words and boundary tables}\label{sec:offset-data}
We encode each zero-sum offset vector by the following alphabet.  Every listed symbol belongs
to \(\{-2,-1,0,1,2\}^3\) and has coordinate sum zero; consequently the row-sum condition
is built into the encoding rather than checked by a separate search.
\begingroup\scriptsize
\[
\begin{array}{c|ccccccc}
\text{symbol}&0&1&2&3&4&5&6\\\hline
\text{offset}&(-2,0,2)&(-1,-1,2)&(-1,0,1)&(0,-2,2)&(0,-1,1)&(0,0,0)&(0,1,-1)\\
\end{array}
\]
\endgroup
\begingroup\scriptsize
\[
\begin{array}{c|cccccc}
\text{symbol}&7&8&9&A&B&C\\\hline
\text{offset}&(1,-2,1)&(1,-1,0)&(1,0,-1)&(2,-2,0)&(2,-1,-1)&(2,0,-2)\\
\end{array}
\]
\endgroup
For a word indexed by \(0,1,\ldots,23\), the first printed symbol is the value at index \(0\).
The complete bulk data reduce to the following two master words:
\[
\begin{array}{c|c}
W_1&\mathtt{8AA2375331378AA537533437}\\
W_2&\mathtt{4331347AA5344334347AA234}
\end{array}
\]
For $n\equiv9\pmod{24}$ use $(W_1,W_2)$ with phases $(6,12)$.  For
$n\equiv21\pmod{24}$ use phases $(10,13)$ and
\begin{equation}\label{eq:c3-word-shift}
 \widehat W_1[p]=W_1[p+16\bmod24],\qquad
 \widehat W_2[p]=W_2[p+13\bmod24].
\end{equation}
Thus the formerly separate words for the second class are exact rotations of the displayed words;
after the phases are inserted, passage between the two classes is a half-period shift by $12$ in
both bulk layers.
For each boundary window, the symbols are listed in increasing order of the indicated relative
index; no reversal of a printed word is intended.
Here \(L,M,U\) mean the windows at \(1,\kappa,n\), respectively.
\[
\begin{array}{c|c|c|l}
n\bmod48&\text{window}&\text{indices}&\text{offset word}\\\hline
9&L&1:20&\mathtt{AAA57AACA8A5331378AA}\\
9&M&-20:20&\mathtt{2375331378AA537A534378A81344331347AA53443}\\
9&U&-20:0&\mathtt{3347137A53432347AAA80}\\
33&L&1:20&\mathtt{57AB88A647A5331378AA}\\
33&M&-20:20&\mathtt{5375334378AA237A530343433344334347AA23443}\\
33&U&-20:0&\mathtt{33431A7A53478547AAAB8}\\
21&L&1:20&\mathtt{A988AAA8AA75334378AA}\\
21&M&-20:20&\mathtt{4378AA2375331378AAA57AA8A547AA2344331347A}\\
21&U&-20:0&\mathtt{3343437A5343334717570}\\
45&L&1:20&\mathtt{A81347A5337533437833}\\
45&M&-20:20&\mathtt{1378AA5375334378AA2375331347A53344334347A}\\
45&U&-20:0&\mathtt{33431337534344333447A}\\
\end{array}
\]
For an independent support audit, fix one of the $16$ admissible periodic states described
below and a residue \(u\in\{1,\ldots,72\}\).  Write $n_0=e+24\tau$, so that
$n=n_0+288Q$ with $Q\ge0$.  Decoding the two bulk words and the corresponding boundary
table turns every family of values congruent to $u\pmod{72}$ into an integer interval
whose endpoints are affine functions of \(Q\).  For each of the $16\times72$ state--residue
pairs, the displayed boundary word gives the endpoint relation at \(Q=0\), while the period
words give the same endpoint slopes for every \(Q\ge0\).  Hence the intervals remain consecutive,
begin at \(0\), and end at
\(12Q+\lfloor(3n_0-u)/72\rfloor\), exactly as required in
Lemma~\ref{lem:rule}.  This is a finite comparison of printed affine expressions, not a
computer-generated existence assertion.

\subsection{The 24 stored orientation patterns for 48 admissible affine states}\label{sec:orientation-data}
For a row index \(j\), let its category be
\[
\gamma(j)=\begin{cases}
(L,j),&j\le20,\\
(M,j-m),&|j-m|\le12,\\
(K,j-\kappa),&|j-\kappa|\le20,\\
(T,j-n),&j\ge n-20,\\
(B_1,j\bmod24),&j<m-12,\\
(B_2,j\bmod24),&m+12<j<\kappa-20,\\
(B_3,j\bmod24),&\text{otherwise}.
\end{cases}
\]
The cases are applied in the displayed order.  For any admissible parameter using the periodic rule,
put \(e=n\bmod24\) (with the representatives \(e=9,21\) in the periodic range),
\(t=(n-e)/24\), and select the stored row with
\(s=15+((t-15)\bmod6)\).  Coprimality gives
\[
s\in\{16,17,19,20\}\quad(e=9),\qquad
s\in\{15,17,18,20\}\quad(e=21).
\]
This also specifies the ten finite-bridge cases below.  For those bridge values, \(s\) is only
the table-selection residue; the affine \(\tau\)-range starts at \(15\).  For the infinite proof, an admissible affine
group state is \((e,\tau,r)\), where
\(n=e+24(\tau+12Q)\), \(e\in\{9,21\}\), \(15\le\tau\le26\), \(Q\ge0\), and
the retained indices satisfy \(j\equiv r\pmod3\).  The stored row \((e,s,r)\),
\(15\le s\le20\), is used for both \(\tau=s\) and \(\tau=s+6\).  For either value of
\(\tau\), let \(\mathcal A_{e,\tau,r}\) be the union of its categories for \(Q=0,1,\ldots,6\);
the literal category set is the same for the two values represented by a stored row.
The restrictions above leave $8$ stored parameter states and hence $24=8\cdot3$ stored rows; they
cover the $48=16\cdot3$ admissible affine group states.  Here \(r\in\{0,1,2\}\) is the group
index, \(p=j\bmod24\in\{0,\ldots,23\}\), and \(z\in\{1,2,3\}\) is the bulk zone.  The
category-to-label dictionary used by the printed strings is
\[
L\longleftrightarrow\mathtt{('bnd','low',q)},\quad
M\longleftrightarrow\mathtt{('bnd','m',q)},\quad
K\longleftrightarrow\mathtt{('bnd','k',q)},\quad
T\longleftrightarrow\mathtt{('bnd','top',q)}.
\]
The letters \(M\) and \(K\) in the category list are deliberately distinct: here \(M\) denotes
the window centred at the integer \(m\), whereas \(K\) denotes the window centred at \(\kappa\).
This category notation is independent of the \(L,M,U\) labels used for the offset-table display
above; the literal labels in the sign strings are the authoritative ones for decoding.
Thus a switch \(K:q\) refers to the literal key \(\mathtt{('bnd','k',q)}\), while a bulk key
has the form \(\mathtt{('bulk',z,p)}\) with \(z\in\{1,2,3\}\) and \(p\in\{0,\ldots,23\}\).
For encoding,
write the labels literally as \(\mathtt{('bnd','k',q)}\), \(\mathtt{('bnd','low',q)}\),
\(\mathtt{('bnd','m',q)}\), \(\mathtt{('bnd','top',q)}\), and
\(\mathtt{('bulk',z,p)}\), and order these literal strings lexicographically.  Each hexadecimal digit below
expands to four binary digits; retain the first \(\ell\) bits, with \(1=+1\) and \(0=-1\).
The columns \(C,D\) give respectively the signs of \(\mu_j\) and \(|w_j|\); these are sign-string
columns and are unrelated to the column-sum vectors \(C_i\) in the finite-bridge table below.
Both columns use the same ordered list of literal keys, whose length is \(\ell\).  A switch
\(X:q;a\) or \(X:q;b\) changes the unique row in category \((X,q)\) from mode \(c\)
to the displayed mode; ``--'' means that every row uses mode \(c\).
\begin{longtable}{c@{\,}c@{\,}c|r|c|l|l}
\hline
e&s&r&$\ell$&\text{switch}&C&D\\\hline
\endfirsthead
\hline e&s&r&$\ell$&\text{switch}&C&D\\\hline
\endhead
9&16&0&59&--&\texttt{FE7FAF801E1E1E0}&\texttt{BFFFFC481FF51C0}\\
9&16&1&60&\(K:-1;b\)&\texttt{FF3B87F00F0B8F0}&\texttt{FFFEE80404DB8F0}\\
9&16&2&60&--&\texttt{FEDFAF00069F0F0}&\texttt{FF7FF0140AA78F0}\\
9&17&0&59&\(K:-11;a\)&\texttt{7FFE55001E1E1E0}&\texttt{BEFBF0081FF51C0}\\
9&17&1&61&--&\texttt{FDFE07C0054F8780}&\texttt{9FFFFC0202CDC780}\\
9&17&2&59&--&\texttt{F7FD0F800D3E1E0}&\texttt{FAFFF808154F1E0}\\
9&19&0&59&\(K:-11;a\)&\texttt{7FFD3C001E1E1E0}&\texttt{F6F7F0081FF51C0}\\
9&19&1&60&--&\texttt{FFABC7C00A9F0F0}&\texttt{FFFF7A04059B8F0}\\
9&19&2&60&--&\texttt{DDF7AF00069F0F0}&\texttt{FBDFF4040AA78F0}\\
9&20&0&59&--&\texttt{FFEBCF801E1E1E0}&\texttt{DFFFFC081FF51C0}\\
9&20&1&61&\(K:-1;b\)&\texttt{FF3B87F00785C780}&\texttt{FFFFFC00026DC780}\\
9&20&2&59&--&\texttt{FFF8AF200D3E1E0}&\texttt{FF7FF028154F1E0}\\
21&15&0&59&--&\texttt{FFDFF000153E1E0}&\texttt{FFBFFE481FBBFC0}\\
21&15&1&61&--&\texttt{FFF847E0054F8780}&\texttt{DFF7FBA202CDC780}\\
21&15&2&59&--&\texttt{BFEF6F800D3E1E0}&\texttt{FFEFFF58154F1E0}\\
21&17&0&59&--&\texttt{EF7BE800153E1E0}&\texttt{FBFFF8081FBBFC0}\\
21&17&1&60&--&\texttt{FEFC87C00A9F0F0}&\texttt{BF7FFE84059B8F0}\\
21&17&2&60&--&\texttt{7FF76F80069F0F0}&\texttt{FFF7FF440AA78F0}\\
21&18&0&59&--&\texttt{FFF52E00153E1E0}&\texttt{BFDFF4081FBBFC0}\\
21&18&1&61&\(K:-10;a\)&\texttt{BFAD87E0054F8780}&\texttt{6F77F80202CDC780}\\
21&18&2&59&\(K:-12;a\)&\texttt{FDDF8F800D3E1E0}&\texttt{EFCFF008154F1E0}\\
21&20&0&59&--&\texttt{FFEC5800153E1E0}&\texttt{7DFFF0081FBBFC0}\\
21&20&1&60&\(K:-10;a\)&\texttt{FEAF17C00A9F0F0}&\texttt{FFAEF804059B8F0}\\
21&20&2&60&\(K:-12;a\)&\texttt{7DFE4F80069F0F0}&\texttt{EFAFF8040AA78F0}\\
\hline
\end{longtable}
For direct checking, substitute the offsets from Section~\ref{sec:offset-data} in the
closed lift formula.  For each admissible stored row, for the two values \(\tau=s,s+6\), and in each category,
the sums of \(\mu_j\) and \(|w_j|\) are polynomials in \(Q\) of degree at most two.  The two decoded
sign rows make their signed totals zero at \(Q=0,1,\ldots,6\) for both values of \(\tau\);
seven zeros therefore force both polynomials to vanish identically in all \(48\) admissible states.
This is a finite arithmetic verification, and no search or program is part of the logical implication.
\subsection{The ten-case finite bridge}\label{sec:bridge-data}
The periodic support rule is used from \(m\ge123\) onward.  The following ten admissible values
of \(m\) bridge the remaining interval \(67\le m\le119\).  Let \(S_m\) be the set of absolute
values produced and let
\(C_i\) be the three-component column-sum vector of the \(i\)-th array.  Direct substitution
in the printed rules gives the complete audit table below.  Here
\(C_i=(C_{i,1},C_{i,2},C_{i,3})\).  Since every row sum is zero by construction,
\(S_m\subseteq[1,9m]\); the equalities \(\min S_m=1\), \(\max S_m=9m\), and
\(|S_m|=9m\) show that \(S_m\) consists of \(9m\) distinct integers in an interval of
length \(9m\), and hence \(S_m=[1,9m]\).  The last column gives \(C_i=\mathbf0\) for
each \(i\), which is exactly the vanishing of all nine column sums.
\[
\begin{array}{c|c|c|c|c|c}
m&n&n\bmod48&s&(\min S_m,\max S_m,|S_m|)&(C_0,C_1,C_2)\\\hline
67&201&9&20&(1,603,603)&(\mathbf{0},\mathbf{0},\mathbf{0})\\
71&213&21&20&(1,639,639)&(\mathbf{0},\mathbf{0},\mathbf{0})\\
79&237&45&15&(1,711,711)&(\mathbf{0},\mathbf{0},\mathbf{0})\\
83&249&9&16&(1,747,747)&(\mathbf{0},\mathbf{0},\mathbf{0})\\
91&273&33&17&(1,819,819)&(\mathbf{0},\mathbf{0},\mathbf{0})\\
95&285&45&17&(1,855,855)&(\mathbf{0},\mathbf{0},\mathbf{0})\\
103&309&21&18&(1,927,927)&(\mathbf{0},\mathbf{0},\mathbf{0})\\
107&321&33&19&(1,963,963)&(\mathbf{0},\mathbf{0},\mathbf{0})\\
115&345&9&20&(1,1035,1035)&(\mathbf{0},\mathbf{0},\mathbf{0})\\
119&357&21&20&(1,1071,1071)&(\mathbf{0},\mathbf{0},\mathbf{0})\\
\end{array}
\]


\section{Three exceptional supports and six preperiod orientations}\label{sec:small-data}
The nine parameters below are recorded in a compact form that uniquely reconstructs every signed
entry.  Only $m=11,19,23$ require exceptional support words.  The six values
$m=31,35,43,47,55,59$ already use the uniform support rule, but precede the disjoint-window range of
the periodic orientation theorem and therefore retain separate orientation certificates.  Three
support words and \(27\) balance records suffice in place of listing all \(969\) rows.

Put \(n=3m\), and let \(a_{i,j}\) be the Kotzig entry defined in
Section~\ref{sec:construction}.  For an offset \(\delta_j=(d_{1j},d_{2j},d_{3j})\), set
\[
T_j=\operatorname{sort}\{\,|3a_{-1,j}+d_{1j}|,\ |3j+d_{2j}|,
|3a_{1,j}+d_{3j}|\,\}=\{a_j,b_j,c_j\},
\qquad a_j<b_j<c_j,\quad a_j+b_j=c_j.
\]
For \(m=31,35,43,47,55,59\), the offsets are exactly those given by the rule in
Section~\ref{sec:offset-data}.
Direct substitution gives a partition of \([1,9m]\) in all six cases.
For the three smaller parameters, use the following radius-two alphabet and words.
\begingroup\scriptsize
\[
\begin{array}{c|ccccccc}
\text{symbol}&0&1&2&3&4&5&6\\\hline
\delta&(-2,0,2)&(-2,1,1)&(-2,2,0)&(-1,-1,2)&(-1,0,1)&(-1,1,0)&(-1,2,-1)\\
\end{array}
\]
\endgroup
\begingroup\scriptsize
\[
\begin{array}{c|ccccccc}
\text{symbol}&7&8&9&A&B&C&D\\\hline
\delta&(0,-2,2)&(0,-1,1)&(0,0,0)&(0,1,-1)&(0,2,-2)&(1,-2,1)&(1,-1,0)\\
\end{array}
\]
\endgroup
\begingroup\scriptsize
\[
\begin{array}{c|ccccc}
\text{symbol}&E&F&G&H&I\\\hline
\delta&(1,0,-1)&(1,1,-2)&(2,-2,0)&(2,-1,-1)&(2,0,-2)\\
\end{array}
\]
\endgroup
The first symbol is \(\delta_1\), and the word is read from left to right:
\[
\begin{array}{c|l}
m&\text{offset word}\\\hline
11&\mathtt{DDGIGGG988478C9C3CDDGGGG47CD97783}\\
19&\mathtt{DD973CDGGDG9CGGGGD2GGGDGG98787873778878887778877CDGHCG983}\\
23&\mathtt{GGGGCDGG97878C39CDDG9847881887778C780787C9778C77CDGEDGGDG97CGDGDGGGD0}\\
\end{array}
\]
It remains to orient the triples.  For \(r=0,1,2\), order the group as
\(G_r=(T_{3t+r+1})_{t=0}^{m-1}\).  Write \((a_t,b_t,c_t)\) for the ordered
entries of its \(t\)-th triple.  The default mode is \(c\); an entry \(t:a\) or
\(t:b\) in the switch column changes only row \(t\) to that mode.  For a row in mode
\(c,a,b\), respectively, define
\[
(\mu_t,\nu_t,x_t)=
(c_t,b_t-a_t,c_t),\quad(a_t,b_t+c_t,a_t),\quad(b_t,a_t+c_t,b_t).
\]
Expand each hexadecimal word into binary, retain its first \(m\) bits, and read
\(1\) as \(+1\) and \(0\) as \(-1\).  Let the words \(E,H\) give the signs
\(\varepsilon_t,\eta_t\), respectively.  The reconstructed row is
\begin{equation}\label{eq:exceptional-row}
R_{r,t}=\left(\frac{\varepsilon_t\mu_t+\eta_t\nu_t}{2},
\frac{\varepsilon_t\mu_t-\eta_t\nu_t}{2},-\varepsilon_t x_t\right).
\end{equation}
\begingroup\scriptsize
\setlength{\LTleft}{0pt}\setlength{\LTright}{0pt}
\begin{longtable}{@{\extracolsep{\fill}}c@{\quad}c@{\quad}c@{\quad}l@{\quad}l@{}}
\hline
\(m\)&\(r\)&\text{switch}&\(E\)&\(H\)\\\hline
\endfirsthead
\hline \(m\)&\(r\)&\text{switch}&\(E\)&\(H\)\\\hline
\endhead
11&0&\(0:b\)&\texttt{256}&\texttt{686}\\
11&1&--&\texttt{156}&\texttt{3C6}\\
11&2&\(0:a\)&\texttt{386}&\texttt{0E6}\\
19&0&\(0:a\)&\texttt{A027E}&\texttt{0225E}\\
19&1&\(0:b\)&\texttt{B013E}&\texttt{0003E}\\
19&2&\(0:a\)&\texttt{0827E}&\texttt{0E03E}\\
23&0&\(0:b\)&\texttt{8005FE}&\texttt{00807E}\\
23&1&\(1:a\)&\texttt{0C40FE}&\texttt{00107E}\\
23&2&--&\texttt{1011FE}&\texttt{0100BE}\\
31&0&--&\texttt{B00017FE}&\texttt{003C10FE}\\
31&1&--&\texttt{C8000FFE}&\texttt{001410FE}\\
31&2&--&\texttt{00003FFE}&\texttt{00B010FE}\\
35&0&\(0:b\)&\texttt{100013FFE}&\texttt{0008811FE}\\
35&1&--&\texttt{080023FFE}&\texttt{0018003FE}\\
35&2&--&\texttt{10000BFFE}&\texttt{0040009FE}\\
43&0&--&\texttt{0000007FFFE}&\texttt{000210017FE}\\
43&1&--&\texttt{9200001FFFE}&\texttt{00070000FFE}\\
43&2&\(0:a\)&\texttt{8040041FFFE}&\texttt{00080800FFE}\\
47&0&--&\texttt{04000027FFFE}&\texttt{000320001FFE}\\
47&1&--&\texttt{E0100003FFFE}&\texttt{000240001FFE}\\
47&2&--&\texttt{00800047FFFE}&\texttt{000180001FFE}\\
55&0&\(0:b\)&\texttt{902000003FFFFE}&\texttt{00004080007FFE}\\
55&1&\(0:a\)&\texttt{004000007FFFFE}&\texttt{00054080007FFE}\\
55&2&--&\texttt{C00400003FFFFE}&\texttt{00002800007FFE}\\
59&0&--&\texttt{000200080FFFFFE}&\texttt{00001420000FFFE}\\
59&1&--&\texttt{008000001FFFFFE}&\texttt{00005100000FFFE}\\
59&2&\(0:a\)&\texttt{828000000FFFFFE}&\texttt{00005001000FFFE}\\
\hline
\end{longtable}
\endgroup
For every row of the table, direct decoding gives
\[
\sum_{t=0}^{m-1}\varepsilon_t\mu_t=0,
\qquad \sum_{t=0}^{m-1}\eta_t\nu_t=0.
\]
Hence (\ref{eq:exceptional-row}) has zero row sums and zero column sums.
The support words and the periodic offset rule give the absolute-value set \([1,9m]\)
in every case.  Thus the displayed words and table are complete, independently checkable
constructions for all nine preperiod values.


\section{The complete \texorpdfstring{$c=4$}{c=4} slice}\label{sec:c4full}

The even parameter $n=4m$ requires period $48$ and four parameter classes, but no new conceptual
ingredient. For this slice, define $a_{i,j}$ by the same Kotzig formula~\eqref{eq:kotzig}, now with
$n=4m$, $r=2n+1$, $i\in\{-1,0,1\}$ and $j\in[-n,n]$, and put $\kappa=n/2$.

\begin{lemma}[even-$n$ Kotzig array]\label{lem:even-kotzig}
For even $n$, formula~\eqref{eq:kotzig} is integer-valued and defines a $3\times(2n+1)$
array whose three rows are permutations of $[-n,n]$, whose columns sum to zero, and which satisfies
$a_{-i,-j}=-a_{i,j}$.
\end{lemma}
\begin{proof}
For $i=0$ one has $a_{0,j}=j$.  Write $j=2q$ or $j=2q+1$.  When $j=2q$, the entries in
the rows $i=1$ and $i=-1$ are respectively $-n/2-q$ and $n/2-q$; when $j=2q+1$, they are
$n/2-q$ and $-n/2-q-1$.  As $q$ ranges over the permitted integers, each of the two noncentral
rows therefore contains every element of $[-n,n]$ exactly once.  The two first terms in
~\eqref{eq:kotzig} cancel between $i=1$ and $i=-1$, while the coefficients of $j$ in the three
rows sum to $-1/2+1-1/2=0$.  This proves the column sums and the stated central symmetry; it also
shows explicitly why the half-integral terms for odd $j$ cancel.
\end{proof}

We use the corrected lift
\begin{equation}\label{eq:c4lift}
R_j=(3a_{-1,j}+d_{-1}(j),\ 3j+d_0(j),\ 3a_{1,j}+d_1(j)).
\end{equation}
Here $1\le j\le n$, and $d(j)=(d_{-1}(j),d_0(j),d_1(j))\in\{-2,-1,0,1,2\}^3$
has coordinate sum zero.
For $\rho=m\bmod12\in\{1,5,7,11\}$, Section~\ref{sec:c4-support-data} prints the two period-$48$
words, the affine phase law and both boundary-table sets. Group rows by $j-1\bmod4$.

\begin{theorem}[the full $c=4$ slice]\label{thm:c4full}
For every odd $m\ge5$ with $\gcd(m,3)=1$, an $\IHS(m,3;4)$ exists.
\end{theorem}
\begin{proof}
Every row in~\eqref{eq:c4lift} sums to zero: the three Kotzig entries in a column sum to zero,
and every displayed offset vector has coordinate sum zero. It remains to verify support and column
balance.
For a periodic class write
$m=m_\rho+12(s_0+12Q)$, where $s_0\in\{0,\ldots,11\}$, $Q\ge0$, and

\[
(m_1,m_5,m_7,m_{11})=(25,29,31,23).
\]
The $4\cdot12=48$ choices of $(\rho,s_0)$ fix all phases, boundary tables and parities. In a bulk
position $j=j_0+48h$, every
absolute coordinate of~\eqref{eq:c4lift} is affine in $(Q,h)$, with $h$-coefficient $\pm72$ or
$\pm144$. Splitting the former by layer parity gives progressions of difference $144$. Direct
substitution of the printed words and tables makes the quotient intervals adjacent in every residue
modulo $144$, from $1$ through $3n$; hence the support is exact. The $Q=0$ endpoint values are given
by the displayed boundary strings, and the fixed affine slopes propagate the adjacency to every
$Q\ge0$. The same finite substitution verifies that the sorted absolute values in every row satisfy
$a+b=c$. The resulting $48\times144$ affine
endpoint comparisons are induced by the compressed finite support certificate printed in Section~\ref{sec:c4-data};
no search output is used in this implication.

For balance, Section~\ref{sec:c4-balance-data} prints one mode-and-two-sign row for each of the
$4\cdot12\cdot4=192$ group states. On the stated periodic range, each signed category total is a
polynomial in $Q$ of degree at most two. The printed signs give zero at $Q=1,\ldots,6$, while $Q=0$
is checked separately, so both independent totals vanish identically. Lemma~\ref{lem:exact} converts them into zero column sums. The periodic
ranges start at $m=25,29,31,23$ in the four residue classes. The only smaller admissible values are
$5,7,11,13,17,19$, and Section~\ref{sec:c4-small-data} gives complete compact constructions for all
six. This covers every admissible $m$.
\end{proof}

\begin{center}
\emph{Example} ($m=7$). The compact data in Section~\ref{sec:c4-small-data} decode to
{\footnotesize
\[
\setlength{\arraycolsep}{2.5pt}
\left[\begin{array}{rrr}
-1&-42&43\\-13&-38&51\\-30&-25&55\\-37&63&-26\\-67&49&18\\73&-11&-62\\75&4&-79
\end{array}\right]\quad
\left[\begin{array}{rrr}
-39&-6&45\\-16&-34&50\\29&-57&28\\61&-41&-20\\-52&69&-17\\-64&74&-10\\81&-5&-76
\end{array}\right]
\]
\[
\left[\begin{array}{rrr}
-7&-40&47\\32&21&-53\\-33&60&-27\\-22&-44&66\\15&-71&56\\-68&77&-9\\83&-3&-80
\end{array}\right]\quad
\left[\begin{array}{rrr}
12&36&-48\\-31&-23&54\\-24&-35&59\\-19&-46&65\\-14&72&-58\\-8&78&-70\\84&-82&-2
\end{array}\right].
\]
}
\end{center}

Every row and column in this example sums to zero, and the four arrays have absolute-value supports
whose union is $[1,84]$.

\section{Explicit finite data for the complete \texorpdfstring{$c=4$}{c=4} slice}\label{sec:c4-data}
This section prints a compressed complete finite witness used in Theorem~\ref{thm:c4full}.
The strings determine every offset, mode and sign; the accompanying programs merely
regenerate and check these data and are not a logical premise of the theorem.

\subsection{The four period-48 support rules}\label{sec:c4-support-data}
The zero-sum offset vectors occurring in the periodic rules are encoded as follows.
\begingroup\scriptsize\[
\begin{array}{c|ccccccc}
\text{symbol}&0&1&2&3&4&5&6\\\hline
\delta&(-2,0,2)&(-1,-1,2)&(-1,0,1)&(0,-2,2)&(0,-1,1)&(0,0,0)&(1,-2,1)\\
\end{array}\]\endgroup
\begingroup\scriptsize\[
\begin{array}{c|ccccc}
\text{symbol}&7&8&9&A&B\\\hline
\delta&(1,-1,0)&(1,0,-1)&(2,-2,0)&(2,-1,-1)&(2,0,-2)\\
\end{array}\]\endgroup
For $\rho\in\{1,5,7,11\}$ put $m=m_\rho+12s$.  The phase used at step $s$
is $(\phi_1+16s,\phi_2+16s)$ modulo $48$.  The parameter and bulk words are
listed below; a word is indexed from $0$.
\[\begin{array}{c|c|c}
\rho&m_\rho&(\phi_1,\phi_2)\\\hline
1&25&(8,20)\\
5&29&(16,16)\\
7&31&(16,16)\\
11&23&(0,0)\\
\end{array}\]
\begingroup\scriptsize
\begin{longtable}{c@{\quad}c@{\quad}l}
\hline $\rho$&\text{word}&\text{offset string}\\\hline\endfirsthead
\hline $\rho$&\text{word}&\text{offset string}\\\hline\endhead
1&$W_{1}$&\texttt{5343679953469A695343679953469A695343679953469A69}\\
1&$W_{2}$&\texttt{9B433343654333469B433343654333469B43334365433346}\\
5&$W_{1}$&\texttt{53469A695343679953469A695343679953469A6953436799}\\
5&$W_{2}$&\texttt{33469B433343654333469B433343654333469B4333436543}\\
7&$W_{1}$&\texttt{5343679953469A695343679953469A695343679953469A69}\\
7&$W_{2}$&\texttt{3343654333469B433343654333469B433343654333469B43}\\
11&$W_{1}$&\texttt{5343679953469A695343679953469A695343679953469A69}\\
11&$W_{2}$&\texttt{3343654333469B433343654333469B433343654333469B43}\\
\hline\end{longtable}\endgroup
Put $n=4m$ and $\kappa=n/2$.  In the two bulk zones use
\[
\begin{aligned}
\delta(j)&=W_1[(j-\phi_1-16s)\bmod48]&& (21\le j\le\kappa-21),\\
\delta(j)&=W_2[(j-\phi_2-16s)\bmod48]&& (\kappa+21\le j\le n-21).
\end{aligned}
\]
On $L=[1,20]$, $M=[\kappa-20,\kappa+20]$ and $T=[n-20,n]$ use
the following strings, according to $s\bmod2$; indices increase from left to right.
\begingroup\scriptsize
\begin{longtable}{c@{\,}c@{\,}c@{\quad}c@{\quad}l}
\hline $\rho$&$s\bmod2$&\text{window}&\text{indices}&\text{offset string}\\\hline\endfirsthead
\hline $\rho$&$s\bmod2$&\text{window}&\text{indices}&\text{offset string}\\\hline\endhead
1&0&L&1:20&\texttt{95679979943679953469}\\
1&0&M&-20:20&\texttt{9953469A695343679997779A653343654333469B4}\\
1&0&T&-20:0&\texttt{334695436943336513469}\\
1&1&L&1:20&\texttt{56543334613679953469}\\
1&1&M&-20:20&\texttt{6953436799534698343343679533469B433343654}\\
1&1&T&-20:0&\texttt{334695469943336B43236}\\
5&0&L&1:20&\texttt{56543336944679953469}\\
5&0&M&-20:20&\texttt{695343679953469A443343679533469B433343654}\\
5&0&T&-20:0&\texttt{367995433343431634236}\\
5&1&L&1:20&\texttt{5679997977A679953469}\\
5&1&M&-20:20&\texttt{9953469A695343654333469A695343654333469B4}\\
5&1&T&-20:0&\texttt{6979654336975343699A7}\\
7&0&L&1:20&\texttt{503434333469A6953436}\\
7&0&M&-20:20&\texttt{469A6953436799534698343343654333469B43334}\\
7&0&T&-20:0&\texttt{334695433346334603236}\\
7&1&L&1:20&\texttt{797334695369A6953436}\\
7&1&M&-20:20&\texttt{43679953469A695343469953679B4333436543334}\\
7&1&T&-20:0&\texttt{331365433344334346997}\\
11&0&L&1:20&\texttt{79779979979973303436}\\
11&0&M&-20:20&\texttt{469A695343675333469A699797754333469B43334}\\
11&0&T&-20:0&\texttt{334695436943336546997}\\
11&1&L&1:20&\texttt{797973433469A6953436}\\
11&1&M&-20:20&\texttt{43679953469A695343469995469B4333436543334}\\
11&1&T&-20:0&\texttt{334695436943369546997}\\
\hline\end{longtable}\endgroup
For a finite support check, fix $\rho$ and write $s=s_0+12Q$.
There are $4\cdot12=48$ states.  In a bulk word write $j=j_0+48h$.
After fixing the sign of each lifted coordinate, every absolute value is affine in $(Q,h)$
with $h$-coefficient $\pm72$ or $\pm144$.  Splitting the $72$-families by $h\bmod2$
gives progressions of difference $144$.  For each residue modulo $144$, substitution of
the printed words and boundary strings gives adjacent quotient intervals beginning at $0$
and ending at $12Q+\lfloor(3n_0-u)/144\rfloor$.  Hence the lifted absolute values
are exactly $[1,3n]$ in all $48$ states.
\subsection{The 192 balance patterns}\label{sec:c4-balance-data}
Group the triples by $j-1\equiv r\pmod4$.  A row category is
$L_q,M_q,T_q$ in the three boundary windows, and $B_{i,p,e}$ in bulk word $i$,
where $p=(j-\phi_i-16s)\bmod48$, $j_0$ is the representative in $[1,48]$
with that position, and $e=((j-j_0)/48)\bmod2$.  Order attained categories as
$L,M,T,B_1,B_2$, numerically within each family.  Expand each hexadecimal string
to binary, retain the first $\ell$ bits, and read $1$ as $+1$ and $0$ as $-1$.
The columns $C,D$ are the signs of $\mu_j$ and $|w_j|$ from Lemma~\ref{lem:exact}.
A displayed switch $L_q:a$ or $L_q:b$ changes that unique row from the default mode $c$.
For $m=m_\rho+12(s_0+12Q)$ use the row $(\rho,s_0,r)$ below.
\begingroup\scriptsize
\setlength{\LTleft}{0pt}\setlength{\LTright}{0pt}
\begin{longtable}{@{\extracolsep{\fill}}c@{\,}c@{\,}c@{\,}r@{\,}c@{\,}l@{\,}l@{}}
\hline $\rho$&$s_0$&$r$&$\ell$&\text{switch}&$C$&$D$\\\hline\endfirsthead
\hline $\rho$&$s_0$&$r$&$\ell$&\text{switch}&$C$&$D$\\\hline\endhead
1&0&0&68&$L_{1}:b$&\texttt{608F7DE232A6B32B1}&\texttt{E007288154B1EA973}\\
1&0&1&69&--&\texttt{D7318A53ECC48DDC48}&\texttt{9BA87478A9E46D6548}\\
1&0&2&68&--&\texttt{AEA626A0200D11FFF}&\texttt{8271B894401C1EBE6}\\
1&0&3&69&--&\texttt{FA994F88367625D918}&\texttt{1B8FD41824151D9F18}\\
1&1&0&68&--&\texttt{70C1C6B2563CF9078}&\texttt{AB8112F0CF2C9C6B8}\\
1&1&1&69&--&\texttt{E788FEB299654AF850}&\texttt{4821D2D9B262678750}\\
1&1&2&68&$L_{3}:a$&\texttt{8EEB8F1F0192DC569}&\texttt{630CF400006A77567}\\
1&1&3&69&--&\texttt{1ACF510D1001F2FEE8}&\texttt{A3071B75906B89A670}\\
1&2&0&68&--&\texttt{B5157404411AF3ED7}&\texttt{331B35435C73610FD}\\
1&2&1&69&$L_{2}:a$&\texttt{A79E82AD2CCB4B52C8}&\texttt{352A48C12481AF9570}\\
1&2&2&68&--&\texttt{9D3FAA2E4B3BA49B2}&\texttt{56C8C989F9211D9AE}\\
1&2&3&69&--&\texttt{AEDF1A8EFB075012F8}&\texttt{580D551E872CA057D8}\\
1&3&0&68&$L_{1}:b$&\texttt{75E162BEE90658F70}&\texttt{4010F30B0203E29BE}\\
1&3&1&69&$L_{6}:a$&\texttt{7122C0F168ED229EC8}&\texttt{CFE356D1F20EBF0850}\\
1&3&2&68&$L_{3}:a$&\texttt{CA6D8F2F0192DC569}&\texttt{42A472D4606B893A7}\\
1&3&3&69&--&\texttt{E8ABEDD9099E764660}&\texttt{BC8A8540398D8AAD98}\\
1&4&0&68&$L_{1}:b$&\texttt{99C5B4AA4BDAA3719}&\texttt{47E7C5671CA66E784}\\
1&4&1&69&--&\texttt{23AFC0AC0805FAF178}&\texttt{72A2690084E2B7EA28}\\
1&4&2&68&--&\texttt{F93C5091028D9FBBA}&\texttt{0E615D2043A1B2E9D}\\
1&4&3&69&--&\texttt{856DC0204893FBD5E0}&\texttt{0801B9100A5369EA58}\\
1&5&0&68&--&\texttt{8D294F1851D8FB891}&\texttt{0117F0E3324D642D7}\\
1&5&1&69&--&\texttt{413468210A17D36778}&\texttt{8C803A62C6257C09D8}\\
1&5&2&68&$L_{3}:a$&\texttt{0F14E00106677FBCA}&\texttt{1105E650FAC65F0A3}\\
1&5&3&69&--&\texttt{A38EA040181C7DBBD8}&\texttt{D0887410504DF0D4D8}\\
1&6&0&68&--&\texttt{60B8C3C9E435EA781}&\texttt{72B82547D0BD0A5CD}\\
1&6&1&69&$L_{2}:a$&\texttt{FDB5057053B4B4E928}&\texttt{826A6090680CCDEB28}\\
1&6&2&68&--&\texttt{291B981A6F33019DF}&\texttt{C040E8B2100D78E33}\\
1&6&3&69&--&\texttt{EDD7204BC7CE92CD50}&\texttt{840C3898D0039B36B0}\\
1&7&0&68&$L_{1}:b$&\texttt{8096D0344105BF6FA}&\texttt{251911ECA02566B87}\\
1&7&1&69&$L_{6}:a$&\texttt{A63BE02035067FAE70}&\texttt{4884BC0810CDF66530}\\
1&7&2&68&$L_{3}:a$&\texttt{093D8F1F0192DC569}&\texttt{0B16E05440C5755EA}\\
1&7&3&69&--&\texttt{2EAA8E7E58216C6670}&\texttt{041D8B002002DB78E8}\\
1&8&0&68&$L_{1}:b$&\texttt{496D966817E7706E4}&\texttt{8F871000C203C3F7C}\\
1&8&1&69&--&\texttt{27156C1DAE90D52BC8}&\texttt{BCE4D2F7AFF5266A10}\\
1&8&2&68&--&\texttt{826741083087F57D6}&\texttt{DA8431088987B06F5}\\
1&8&3&69&--&\texttt{B7B226A4EA4F41F1C0}&\texttt{63505818324758A1F8}\\
1&9&0&68&--&\texttt{968C8EBA6213B20FC}&\texttt{AA4DCC84392BCD9D0}\\
1&9&1&69&--&\texttt{B5A6461B05DDA18738}&\texttt{56B9C937D44EBE4E00}\\
1&9&2&68&$L_{3}:a$&\texttt{C9ED8F2F0192DC569}&\texttt{AE04725AB40528AF7}\\
1&9&3&69&--&\texttt{ABAB7E012102FBF2E8}&\texttt{D72182B598E4546B58}\\
1&10&0&68&--&\texttt{350D6281021FB1F3B}&\texttt{B63A2017057B321EE}\\
1&10&1&69&$L_{2}:a$&\texttt{C8BA03554AD532B4B0}&\texttt{0F3FC869E53A2EEC50}\\
1&10&2&68&--&\texttt{2E41C7DE304655D92}&\texttt{84E52400E829A15FB}\\
1&10&3&69&--&\texttt{B5F008200B85FBB730}&\texttt{18836432C84CA783F0}\\
1&11&0&68&$L_{1}:b$&\texttt{209B66B51D8629763}&\texttt{8404D08A0005FBC4E}\\
1&11&1&69&$L_{6}:a$&\texttt{59C288E267AF20F0B8}&\texttt{44A960D0104B59AAD8}\\
1&11&2&68&$L_{3}:a$&\texttt{1BAD8F2F0192DC569}&\texttt{D950CBFDFF7BA2428}\\
1&11&3&69&--&\texttt{AAA8A080920FDBAE98}&\texttt{9D325B78C9C6D8A470}\\
5&0&0&68&--&\texttt{0AADB900290736DFB}&\texttt{4CA3D4CD3719432BB}\\
5&0&1&69&$L_{6}:a$&\texttt{158F5DEDF9FA780440}&\texttt{8542F11102A5057DF0}\\
5&0&2&68&--&\texttt{CF2C6DB69C027BAD0}&\texttt{2681B8520677295CE}\\
5&0&3&69&--&\texttt{CEE388F25671797680}&\texttt{80007C640049C6F4D8}\\
5&1&0&68&$L_{1}:b$&\texttt{EB08CD97581486E3B}&\texttt{FFFFD5FB88155E545}\\
5&1&1&69&$L_{6}:a$&\texttt{9B04F4C399AB0139F8}&\texttt{C3D9671B432E35EC40}\\
5&1&2&68&$L_{3}:a$&\texttt{D3636CFD7EF4F5000}&\texttt{9424703478AE0956F}\\
5&1&3&69&--&\texttt{E747AE1641FAE5F280}&\texttt{AF852C14742ED4A2D8}\\
5&2&0&68&$L_{1}:b$&\texttt{80BC529CEE23817D9}&\texttt{A6A0B0B81C397234F}\\
5&2&1&69&--&\texttt{C026B8C4010FB8EDB8}&\texttt{8914D2914087396998}\\
5&2&2&68&--&\texttt{CE07B204A10E3ACFF}&\texttt{84EE2810588F605DD}\\
5&2&3&69&--&\texttt{A3D485BBFFD7C08800}&\texttt{DC0F55D8AA2BE038D8}\\
5&3&0&68&--&\texttt{430E91A75D1BC868B}&\texttt{54D98A93403B0EC2F}\\
5&3&1&69&--&\texttt{BD41A53334AD639170}&\texttt{D8EA76EDFE38CB32C0}\\
5&3&2&68&$L_{3}:a$&\texttt{D6E1EEC6A8AD1C8D5}&\texttt{9C48BC8042167E553}\\
5&3&3&69&--&\texttt{ECC3943CD661F15468}&\texttt{244BD7B0CE456CD2A0}\\
5&4&0&68&--&\texttt{FC718C00045BD0BFF}&\texttt{D4545F18490956FD0}\\
5&4&1&69&$L_{6}:a$&\texttt{5876B2F9F7FC340260}&\texttt{4226B0D100C1D96B58}\\
5&4&2&68&--&\texttt{C70B3A15000BFD96E}&\texttt{DA7921800CAF8AE5A}\\
5&4&3&69&--&\texttt{EAA28C989D9B29B138}&\texttt{F3AC25F0408A6AB478}\\
5&5&0&68&$L_{1}:b$&\texttt{7BAC8A10410AFC6EF}&\texttt{D0613F80960F4933A}\\
5&5&1&69&$L_{6}:a$&\texttt{800677DFF478741440}&\texttt{283BAADDE40849F2E8}\\
5&5&2&68&$L_{3}:a$&\texttt{777B5B03E5CD97541}&\texttt{5D461C3BFE1984AAB}\\
5&5&3&69&--&\texttt{595C2A66D4774403F8}&\texttt{1BB1EDEB8FA72A84C8}\\
5&6&0&68&$L_{1}:b$&\texttt{1B15A958BB4082F7D}&\texttt{179A3C230E4CDB631}\\
5&6&1&69&--&\texttt{8B29CB480401E6FF70}&\texttt{7F032C9C876D6BE180}\\
5&6&2&68&--&\texttt{3203F050190B1FEED}&\texttt{6805A0482A8CB5BB1}\\
5&6&3&69&--&\texttt{AAA2AD646D3093B9A8}&\texttt{4002DAF9826AD1CA68}\\
5&7&0&68&--&\texttt{F2D868A8EAEA69A4E}&\texttt{A011B17A610E4EE16}\\
5&7&1&69&--&\texttt{0BF99B81973A3B9838}&\texttt{A50AAA081E0174BE70}\\
5&7&2&68&$L_{3}:a$&\texttt{8BC83E81C3D1B2B1E}&\texttt{A8FD8FB73D7D80C87}\\
5&7&3&69&--&\texttt{B160FE8477648F3E40}&\texttt{073164809503781778}\\
5&8&0&68&--&\texttt{193D5022205BE97BE}&\texttt{F8BD440DDB3998E9A}\\
5&8&1&69&$L_{6}:a$&\texttt{867BB5FDAF7C724040}&\texttt{608188831127C26EC8}\\
5&8&2&68&--&\texttt{4A8B55597C4A1D473}&\texttt{6C62B347E07BF818C}\\
5&8&3&69&--&\texttt{8E622BBDFDDE380220}&\texttt{3CE4B4BF05489CE570}\\
5&9&0&68&$L_{1}:b$&\texttt{63A2DC049009E7F5B}&\texttt{5C8436BD02D9BC866}\\
5&9&1&69&$L_{6}:a$&\texttt{3CBCAF4EBFFC620620}&\texttt{38111880B4875517B0}\\
5&9&2&68&$L_{3}:a$&\texttt{DBC41BBFEFC4C0C0C}&\texttt{835AE244A80519EFC}\\
5&9&3&69&--&\texttt{4FE908A4001FFF2E90}&\texttt{731C7B208CE30C8FF0}\\
5&10&0&68&$L_{1}:b$&\texttt{2D39C194040F6FCF8}&\texttt{D03B3284DBB90D2AF}\\
5&10&1&69&--&\texttt{EB492141800B9EBE78}&\texttt{CE09DAF78834A75E40}\\
5&10&2&68&--&\texttt{C54675E91D8F18A4D}&\texttt{18925540144D15D8F}\\
5&10&3&69&--&\texttt{A2E646400067744FF8}&\texttt{A3449C080026E7F0E0}\\
5&11&0&68&--&\texttt{425A9BB307125565E}&\texttt{6BFA1088F8AB9CCAC}\\
5&11&1&69&--&\texttt{F8189840046DB77FC0}&\texttt{1946987331F1F80DC8}\\
5&11&2&68&$L_{3}:a$&\texttt{5E255EFF7CB0F040C}&\texttt{972256692249293DB}\\
5&11&3&69&--&\texttt{FF04DBF7DF76994000}&\texttt{996A2E84D2DC48C5F8}\\
7&0&0&68&$L_{1}:b$&\texttt{6F328C96C338D143F}&\texttt{BABC261DA2D2BE591}\\
7&0&1&69&$L_{2}:a$&\texttt{7440F0379AC90D99E8}&\texttt{8609294300C49CE978}\\
7&0&2&68&$L_{3}:a$&\texttt{726757FBD3F3F2000}&\texttt{9B80B812448CA86BF}\\
7&0&3&69&--&\texttt{2BCA99C00887BFDD40}&\texttt{C566195B51E1B1DE20}\\
7&1&0&68&--&\texttt{981DB0161805E9FEB}&\texttt{57CA9004B226D07EE}\\
7&1&1&69&$L_{2}:a$&\texttt{2D142CE92D347625B0}&\texttt{B521503B6BC7C32930}\\
7&1&2&68&$L_{3}:a$&\texttt{F5EFD58AC5D5DA85C}&\texttt{FFFFB57DA1455A754}\\
7&1&3&69&$L_{20}:a$&\texttt{D8D3B48CDD630C73C8}&\texttt{4C62C04001259D8F38}\\
7&2&0&68&--&\texttt{C40B347E1871C355E}&\texttt{F3794C9A7000C3E6F}\\
7&2&1&69&--&\texttt{4901F24040A5ECBBD8}&\texttt{CA42F29568EF561518}\\
7&2&2&68&$L_{3}:a$&\texttt{D43FCDFEBFA8E2202}&\texttt{0C157098C5CC798F4}\\
7&2&3&69&--&\texttt{03D739500066FCEF50}&\texttt{84007805A123D9C2B8}\\
7&3&0&68&$L_{5}:b$&\texttt{164FA1801413FA6BF}&\texttt{AEFE64C5306F103E7}\\
7&3&1&69&--&\texttt{7C63520112159DCFF0}&\texttt{3C162B810441ADC758}\\
7&3&2&68&$L_{3}:a$&\texttt{31C4E5D18DA954DAA}&\texttt{DC4B8C40E0E5B46BA}\\
7&3&3&69&$L_{20}:a$&\texttt{C10F48B7CA6199BA90}&\texttt{8BF17530EA6F603338}\\
7&4&0&68&$L_{1}:b$&\texttt{16BE1E01F2EC455B3}&\texttt{92F72C19740D49967}\\
7&4&1&69&$L_{2}:a$&\texttt{BF9F3860842752FAF8}&\texttt{FFFFE2AD54ADD23550}\\
7&4&2&68&$L_{3}:a$&\texttt{A819DAFFFEA067090}&\texttt{5A1766D538A90DE35}\\
7&4&3&69&--&\texttt{F9115431AEC8EE82B8}&\texttt{6B1547946A30232F78}\\
7&5&0&68&--&\texttt{5007F820118F5ACFE}&\texttt{DB499B98D2613C4F5}\\
7&5&1&69&$L_{2}:a$&\texttt{6EC5AB4B4AAACB4B28}&\texttt{AC00384483860CDF98}\\
7&5&2&68&$L_{3}:a$&\texttt{7188B1805083B4EFF}&\texttt{3059D02160C895EDE}\\
7&5&3&69&$L_{20}:a$&\texttt{44FD52216005BBE6E8}&\texttt{84F1ACB03E023BF148}\\
7&6&0&68&--&\texttt{E27443A5CCA59A365}&\texttt{10093555D52A44FB4}\\
7&6&1&69&--&\texttt{A87CB04A84047FB770}&\texttt{5485D88C005D073AF8}\\
7&6&2&68&$L_{3}:a$&\texttt{AF16729A8BEE32527}&\texttt{D144D840883AB3FA8}\\
7&6&3&69&--&\texttt{6074A3009404F37FA8}&\texttt{D821208AE9E929CEA8}\\
7&7&0&68&$L_{5}:b$&\texttt{4A357014C204B9FBF}&\texttt{5237220DDE3DC92A6}\\
7&7&1&69&--&\texttt{F4171004A02DF6D798}&\texttt{B8458E00858529F670}\\
7&7&2&68&$L_{3}:a$&\texttt{C9A99F98C192DC569}&\texttt{55D4BBFFFF63B20C2}\\
7&7&3&69&$L_{20}:a$&\texttt{9D2B4EE5E8456D6118}&\texttt{7CC42C26885E3B4D88}\\
7&8&0&68&$L_{1}:b$&\texttt{6BF089A313DF2982E}&\texttt{050DF050D665AEE13}\\
7&8&1&69&$L_{2}:a$&\texttt{F84358050E07AB5FD0}&\texttt{8830B211580599D750}\\
7&8&2&68&$L_{3}:a$&\texttt{B6079B55E82D4DC70}&\texttt{A103B75F53A43F44C}\\
7&8&3&69&--&\texttt{831C98BD33C10537D8}&\texttt{82CBB834B864FF0C90}\\
7&9&0&68&--&\texttt{DB817308060FED69D}&\texttt{7F420A102092FCE56}\\
7&9&1&69&$L_{2}:a$&\texttt{86756F6B0934AD2AB0}&\texttt{5DA028BAA2F51732A8}\\
7&9&2&68&$L_{3}:a$&\texttt{2B12BFBDEF50D80E0}&\texttt{ED44CCD309736C996}\\
7&9&3&69&$L_{20}:a$&\texttt{44F388AC76B6D4B248}&\texttt{E7D45CB01009A9BD38}\\
7&10&0&68&--&\texttt{CB823D51559B88697}&\texttt{9456CC701D0C97467}\\
7&10&1&69&--&\texttt{233B20010A55E9AFB8}&\texttt{D211983C08857E0DB0}\\
7&10&2&68&$L_{3}:a$&\texttt{29F5B96A7469A47AA}&\texttt{80A13C910E17074BD}\\
7&10&3&69&--&\texttt{3C0BBB3BDFF8029868}&\texttt{572A3CAC130CE899E8}\\
7&11&0&68&$L_{5}:b$&\texttt{922AD572A17D898E6}&\texttt{88ACB71011E6F865A}\\
7&11&1&69&--&\texttt{D946C758D9241D13E8}&\texttt{01FBC4C9FC22D59628}\\
7&11&2&68&$L_{3}:a$&\texttt{458F1D97BFF30A070}&\texttt{A229C0480201E774F}\\
7&11&3&69&$L_{20}:a$&\texttt{0C87E80E0183FEED30}&\texttt{BBCA2F081B4533B298}\\
11&0&0&68&$L_{5}:b$&\texttt{6DB94B8CA9CEA49F0}&\texttt{8B327B10C230F2E97}\\
11&0&1&69&--&\texttt{09D51ABDFBFA068540}&\texttt{3D9A4941C2071DBA48}\\
11&0&2&68&$L_{3}:a$&\texttt{4011F2C2001AAFDF6}&\texttt{5CC3539F8498913ED}\\
11&0&3&69&$L_{4}:a$&\texttt{FFC861FFBFE0A41888}&\texttt{CD013C5D90855E6268}\\
11&1&0&68&$L_{9}:b$&\texttt{4446BA000C8FB6AFA}&\texttt{9BD89BA17344F01F9}\\
11&1&1&69&--&\texttt{DCB960639CDD6BD880}&\texttt{19882FF745FC089770}\\
11&1&2&68&$L_{3}:a$&\texttt{66B7AE7FFCE0FA800}&\texttt{6263A2D6DA424FAE8}\\
11&1&3&69&$L_{20}:a$&\texttt{1F86945CB86B246C78}&\texttt{4D035934FC60156EE8}\\
11&2&0&68&--&\texttt{86587435D35983E35}&\texttt{144D948002036BD3D}\\
11&2&1&69&$L_{2}:a$&\texttt{FDE9B3C52C5BE57108}&\texttt{E62C98051C0F7F5810}\\
11&2&2&68&$L_{3}:a$&\texttt{5CCFC7EDDFE47A800}&\texttt{1AC736D77C63158F4}\\
11&2&3&69&$L_{4}:a$&\texttt{AF1BE5A2CC3EF904C8}&\texttt{E20E5930CBE96151F8}\\
11&3&0&68&--&\texttt{B02DAEDF804ED3449}&\texttt{0C23A1BC4818FB295}\\
11&3&1&69&$L_{2}:a$&\texttt{B1E9CCB4B55534B4D0}&\texttt{D75350A45BF2A2EE30}\\
11&3&2&68&$L_{3}:a$&\texttt{6AD4DB7777F078A20}&\texttt{9127B3F797F09BEA0}\\
11&3&3&69&$L_{20}:a$&\texttt{5256DC8F5A60655ED0}&\texttt{10E93FB12A36CB0698}\\
11&4&0&68&$L_{5}:b$&\texttt{EA4D2F19F08350F5C}&\texttt{3248EA2B304444DFD}\\
11&4&1&69&--&\texttt{24AE391C001377ECE8}&\texttt{B000E8C521046AF8E8}\\
11&4&2&68&$L_{3}:a$&\texttt{029DB5DFEDF4F8001}&\texttt{1814C698A84D338DC}\\
11&4&3&69&$L_{4}:a$&\texttt{43722BFFF3B680C440}&\texttt{4E4A5C10C809BB31B8}\\
11&5&0&68&$L_{9}:b$&\texttt{4ADD8100036FFB237}&\texttt{A1046634002C1D7E6}\\
11&5&1&69&--&\texttt{ECB4E0EB3B10DE89E0}&\texttt{C6A888D760BC69CF80}\\
11&5&2&68&$L_{3}:a$&\texttt{36330FEF9BD472082}&\texttt{E409302C02CEAD63C}\\
11&5&3&69&$L_{20}:a$&\texttt{36EAF0FEA9236092F8}&\texttt{2EC4B485C2650B62F8}\\
11&6&0&68&--&\texttt{831A9B2F12AC50BB3}&\texttt{8155EDA33005BE741}\\
11&6&1&69&$L_{2}:a$&\texttt{6F05E52D2CCCAD2D48}&\texttt{6C855128480BDAB9C0}\\
11&6&2&68&$L_{3}:a$&\texttt{F7B89614E97CA94C7}&\texttt{0D08F0385CCDBE0CC}\\
11&6&3&69&$L_{4}:a$&\texttt{26768044A042FBB7E0}&\texttt{0FFCF575FD81B026F0}\\
11&7&0&68&--&\texttt{BE471F45D0C9C991D}&\texttt{5C60412014AE543DE}\\
11&7&1&69&$L_{2}:a$&\texttt{09BD83316B1D09D568}&\texttt{8C627E6F5E63D31488}\\
11&7&2&68&$L_{3}:a$&\texttt{2B30E7FBF3D8F0140}&\texttt{A77859100C15E5769}\\
11&7&3&69&$L_{20}:a$&\texttt{249CDA010495F5B678}&\texttt{F21849020203E8E5F0}\\
11&8&0&68&$L_{5}:b$&\texttt{499FC9201807E6FCB}&\texttt{13116940ABFBE4919}\\
11&8&1&69&--&\texttt{648E684820A73DB9D8}&\texttt{A54DB64682646DD720}\\
11&8&2&68&$L_{3}:a$&\texttt{0DB96FE76FDC7C000}&\texttt{712390C4C0267874F}\\
11&8&3&69&$L_{4}:a$&\texttt{6507CB7FBB7E480848}&\texttt{6514584542475A71C8}\\
11&9&0&68&$L_{9}:b$&\texttt{319AAB58CE4B2B2D4}&\texttt{27C46C04F21CBBC2A}\\
11&9&1&69&--&\texttt{91A7B0DEA17041DEE8}&\texttt{6006BA01A5028BDF28}\\
11&9&2&68&$L_{3}:a$&\texttt{9F2097CEFFE8C80C2}&\texttt{AE735134589ED0BB8}\\
11&9&3&69&$L_{20}:a$&\texttt{78867495DB81F92B80}&\texttt{E31734B4122F18B728}\\
11&10&0&68&--&\texttt{0A4F5D6D42AB4B0D5}&\texttt{3B68E6140036DE78C}\\
11&10&1&69&$L_{2}:a$&\texttt{201ADB594D52DB5C80}&\texttt{65FE101400BBED2B48}\\
11&10&2&68&$L_{3}:a$&\texttt{55D8ACF7DF74B80C0}&\texttt{E026A440100BC8FE6}\\
11&10&3&69&$L_{4}:a$&\texttt{51487E9E506858B5D8}&\texttt{DB1311200043F7ED00}\\
11&11&0&68&--&\texttt{1A8ED6753A13DA478}&\texttt{141F80CCC5461CEA7}\\
11&11&1&69&$L_{2}:a$&\texttt{508E6C080524F77DE0}&\texttt{D06524830283EF6B00}\\
11&11&2&68&$L_{3}:a$&\texttt{CE588DFE77E4780C0}&\texttt{C6C31716D7B8CD385}\\
11&11&3&69&$L_{20}:a$&\texttt{DC48EEE0BE0AC26B58}&\texttt{BE9EA431033EE087F0}\\
\hline\end{longtable}\endgroup
For a fixed row of this table, every category contribution is eventually quadratic in $Q$.
Direct substitution gives zero for $Q=1,\ldots,6$ in both signed chains, while $Q=0$
is checked separately; hence both totals vanish identically.  These $192=4\cdot12\cdot4$
rows are therefore the complete finite orientation certificate.
\subsection{The six bounded constructions}\label{sec:c4-small-data}
For $m\in\{5,7,11,13,17,19\}$ the following compact data reconstruct every signed entry.
For each such $m$, they produce four groups of $m$ rows, hence four $m\times3$ arrays.
The radius-two offset alphabet is
\begingroup\scriptsize\[
\begin{array}{c|ccccccc}
\text{symbol}&0&1&2&3&4&5&6\\\hline
\delta&(-2,0,2)&(-2,1,1)&(-2,2,0)&(-1,-1,2)&(-1,0,1)&(-1,1,0)&(-1,2,-1)\\
\end{array}\]\endgroup
\begingroup\scriptsize\[
\begin{array}{c|ccccccc}
\text{symbol}&7&8&9&A&B&C&D\\\hline
\delta&(0,-2,2)&(0,-1,1)&(0,0,0)&(0,1,-1)&(0,2,-2)&(1,-2,1)&(1,-1,0)\\
\end{array}\]\endgroup
\begingroup\scriptsize\[
\begin{array}{c|ccccc}
\text{symbol}&E&F&G&H&I\\\hline
\delta&(1,0,-1)&(1,1,-2)&(2,-2,0)&(2,-1,-1)&(2,0,-2)\\
\end{array}\]\endgroup
The word gives $\delta_1,\ldots,\delta_{4m}$ in the corrected Kotzig lift.
\begingroup\scriptsize
\begin{longtable}{c@{\quad}l}
\hline $m$&\text{offset word}\\\hline\endfirsthead\hline $m$&\text{offset word}\\\hline\endhead
5&\texttt{D37CDDDD9C9CGGGGDDG4}\\
7&\texttt{G9C97CEDGD98738CGGDGHCDGIGD0}\\
11&\texttt{G9707478CGGD9CDDDGE9C788777877CGIGDGHCDGIGD0}\\
13&\texttt{GGDGDDGGB778CGDGD77CD978CGDGGDD777IGGGGGD988887CGGD0}\\
17&\texttt{DGDDDDG2CGG87CGDDDA8CDDD988874DDD7C7CDG978777CGGGG94GDG8788CD9C7C973}\\
19&\texttt{D58CG9CDGCGCDGG9CGDGDDGDA4DDGCDGDDDD987CDG907877787C9777777878CGHDDGGGD7H983}\\
\hline\end{longtable}\endgroup
Let $T_j=\{a_j,b_j,c_j\}$, $a_j<b_j<c_j=a_j+b_j$, be the sorted triple
obtained from the word, and order group $r$ as $T_{4t+r+1}$.
For mode $c,a,b$, respectively, put
\[
(\mu_t,\nu_t,x_t)=(c_t,b_t-a_t,c_t),\quad(a_t,b_t+c_t,a_t),\quad(b_t,a_t+c_t,b_t).
\]
Decode $E,H$ as signs $\varepsilon_t,\eta_t$ in row order and set
\begin{equation}\label{eq:c4-small-row}
R_{r,t}=\left((\varepsilon_t\mu_t+\eta_t\nu_t)/2,
(\varepsilon_t\mu_t-\eta_t\nu_t)/2,-\varepsilon_t x_t\right).
\end{equation}
The switch column lists the nondefault modes; ``--'' means all modes are $c$.
\begingroup\scriptsize
\setlength{\LTleft}{0pt}\setlength{\LTright}{0pt}
\begin{longtable}{@{\extracolsep{\fill}}c@{\quad}c@{\quad}l@{\quad}l@{\quad}l@{}}
\hline $m$&$r$&\text{switches}&$E$&$H$\\\hline\endfirsthead
\hline $m$&$r$&\text{switches}&$E$&$H$\\\hline\endhead
5&0&$2:b,4:a$&\texttt{C8}&\texttt{48}\\
5&1&$1:a,2:a,4:b$&\texttt{88}&\texttt{48}\\
5&2&$2:b$&\texttt{18}&\texttt{18}\\
5&3&$1:a,2:a,3:a,4:b$&\texttt{18}&\texttt{18}\\
7&0&$3:a,4:a,5:b$&\texttt{16}&\texttt{C6}\\
7&1&$2:a,3:a,4:a,5:a,6:b$&\texttt{1E}&\texttt{72}\\
7&2&$2:a,4:b,5:a,6:b$&\texttt{66}&\texttt{DA}\\
7&3&$4:b,5:b,6:a$&\texttt{8E}&\texttt{32}\\
11&0&--&\texttt{436}&\texttt{0DA}\\
11&1&--&\texttt{82E}&\texttt{116}\\
11&2&$3:b$&\texttt{60E}&\texttt{2BA}\\
11&3&$2:b$&\texttt{21E}&\texttt{89A}\\
13&0&$0:b$&\texttt{8878}&\texttt{3058}\\
13&1&--&\texttt{48B8}&\texttt{0A38}\\
13&2&$0:a$&\texttt{BA18}&\texttt{0D38}\\
13&3&$2:b$&\texttt{2478}&\texttt{4C38}\\
17&0&--&\texttt{C20F8}&\texttt{14078}\\
17&1&$0:a$&\texttt{81578}&\texttt{34078}\\
17&2&$0:a$&\texttt{80E78}&\texttt{12338}\\
17&3&$0:a$&\texttt{021F8}&\texttt{02478}\\
19&0&$1:b$&\texttt{6185E}&\texttt{0505E}\\
19&1&--&\texttt{D013E}&\texttt{02D8E}\\
19&2&--&\texttt{3603E}&\texttt{8C02E}\\
19&3&--&\texttt{5903E}&\texttt{0511E}\\
\hline\end{longtable}\endgroup
For every displayed row, decoding gives $\sum\varepsilon_t\mu_t=0$ and
$\sum\eta_t\nu_t=0$.  Formula~\eqref{eq:c4-small-row} therefore has zero row and
column sums, while the support words give each absolute value in $[1,12m]$ exactly once.


\section{Concluding remarks}\label{sec:remarks}

The construction separates the width-three problem into three independently checkable parts: a
global support partition, a grouping into arrays, and an orientation certificate for each group. The
centre-complementary symmetry of the Kotzig formula reduces the support check to periodic bulk rules
and finitely many boundary corrections, and it also organises the orientation data. Related
centre-complementary machinery underlies the construction of magic three-dimensional rectangles
in~\cite{Zhou2018JCD}; here only one half of the symmetric array is used, because every absolute
value must occur exactly once.

The same decomposition suggests a route for other values of $c$. In this paper the period-$24$
certificate closes the odd slice $c=3$, while an independent period-$48$ certificate closes the even
slice $c=4$. When $n=mc$ is even, both the hole correction and the orientation categories have
different parity behaviour, so the period-$24$ certificate does not transfer without a new argument.
Establishing certificates for $c\ge5$ remains open here.

\paragraph{Code and data availability.}
Section~\ref{sec:certificates} contains the two master offset words, their exact rotation rule, all
four boundary-table sets, all $24$ admissible orientation patterns, and the ten-case finite bridge.
The $m=7$ construction appears in the main proof, and Section~\ref{sec:small-data} contains complete
compact reconstruction data for three exceptional supports and six preperiod orientations.
For $c=4$, Section~\ref{sec:c4-data} contains the four period-$48$ support rules, the $192$ balance
rows, and the six compact seed constructions. Consequently both fixed-slice theorems can be checked
from the printed manuscript together with the accompanying certificate source package. The public
repository at \url{https://github.com/matheorist/heffter-width3-c34} contains the exact source files,
machine-readable data, and verification suite; its README gives the file inventory, dependencies,
and commands. The deterministic protocol is to decode the offset strings, reconstruct and sort every
lifted row, verify the row sum-triple and support conditions, and then decode the two sign strings and
verify the two signed totals for each group. The programs are retained for reproducibility and
convenience, but none is needed as an oracle or as a logical premise of the result.

\section*{Funding}
\begingroup
\small
\begin{minipage}{0.92\textwidth}
\raggedright
Yong Zhang was supported by the National Natural Science Foundation of China
(Grant No.~11871417) and the 2026 Jiangsu Provincial Teacher Development Research
Project (Grant No.~2026jsfz-b004). Guangzhou Chen was supported by the National
Natural Science Foundation of China (Grant No.~12471245) and the Natural Science
Foundation of Henan Province (Grant No.~262300421846).
\end{minipage}
\endgroup

\begingroup
\footnotesize
\bibliographystyle{plain}
\bibliography{heffter_c3}
\endgroup

\end{document}
